% Generated draft: requires mathematical and source review.
\documentclass[11pt]{article}
\usepackage[margin=25mm]{geometry}
\usepackage{fontspec}
\setmainfont{Latin Modern Roman}
\usepackage{amsmath,amssymb,mathtools,unicode-math}
\setmathfont{Latin Modern Math}
\usepackage{enumitem,needspace,xurl,hyperref}
\setlength{\parindent}{0pt}
\setlength{\parskip}{5pt}
\setlength{\emergencystretch}{3em}
\newcommand{\N}{\mathbb{N}}
\newcommand{\Z}{\mathbb{Z}}
\newcommand{\Q}{\mathbb{Q}}
\newcommand{\R}{\mathbb{R}}
\newcommand{\C}{\mathbb{C}}
\newcommand{\F}{\mathbb{F}}
\newcommand{\A}{\mathbb{A}}
\newcommand{\PP}{\mathbb{P}}
\newcommand{\E}{\mathbb{E}}
\newcommand{\eps}{\varepsilon}
\newcommand{\dd}{\,\mathrm{d}}
\begin{document}
\section*{UnsolvedMath ID: 11000261}
\textbf{Source problem code:} catalogue entry
\subsection{Definitions and mathematical statement}
Let \(\operatorname{BMW}_n\) denote the Birman--Wenzl--Murakami algebra (also ordered Birman--Murakami--Wenzl in the source), in the coefficient-ring, parameter and rank conventions intended by the chapter. The problem asks which constructions and results of ``this paper'' extend to \(\operatorname{BMW}_n\), and specifically proposes generalizing [Big04] to this algebra. The paper's relevant theorems, Section 7 and [Big04] are not reproduced, so no particular homological, braid-group or representation construction is invented as their content.

The source supplies a starting result: John Enyang [Eny04] proved a cellular-algebra structure and used it to define irreducible representations by an approach similar to Section 7. In a standard cellular framework, cell modules and invariant bilinear forms provide candidate simple modules through their nonzero quotients by form radicals, with the precise conclusion depending on the field and specialization. This explains the kind of representation-theoretic input reported, not a claim that it already extends every result or that all specializations are semisimple.

A precise extension must identify the original construction/theorem, define its BMW analogue with the chosen parameter hypotheses, and state the comparison or property that survives. The request is broader than verifying cellularity, which the source already reports. It gives no prescribed parameter regime, field characteristic, generic/specialized case or list of additional theorems. The specific next goal remains the unnamed [Big04] generalization until that reference is supplied.
\subsection{Short English statement}
Determine which results of the omitted paper extend to the BMW algebra, building on Enyang's reported cellularity, with [Big04] as the next specific extension target. The actual constructions and parameter conventions are missing.
\subsection{Sources}
\begin{itemize}
\item Local UnsolvedMath registry, permanent ID \(11000261\), exact source code \texttt{catalogue entry}.
\item Farb, \emph{Problems on Mapping Class Groups and Related Topics}, Question 4, PDF page 322: \url{https://www.math.uchicago.edu/~farb/papers/mcgbook.pdf}. Registry access 2026-07-29; pdf-font-structure extraction.
\item John Enyang [Eny04] is cited for cellularity and its representation application; [Big04] is the proposed next generalization. Full titles, statements and bibliography are absent.
\end{itemize}
\end{document}
